\documentclass{article}
\usepackage[utf8]{inputenc}
\usepackage[spanish]{babel}
\usepackage{amsmath}
\usepackage{amssymb}
\usepackage{enumitem}

\title{Ejercicios Funciones Complejas}
\author{}
\date{}

\begin{document}

\maketitle

\section*{Ejercicio 4 (1.1.8 Vol. 2)}

\textbf{Encuentre las derivadas:}

\subsection*{a) $f(z) = \frac{z - i}{z + i}$, encontrar $f'(i)$}

Aplicando la regla del cociente:
\begin{align*}
f'(z) &= \frac{(z + i)(1) - (z - i)(1)}{(z + i)^2} \\
      &= \frac{z + i - z + i}{(z + i)^2} \\
      &= \frac{2i}{(z + i)^2}
\end{align*}

Evaluando en $z = i$:
\begin{align*}
f'(i) &= \frac{2i}{(i + i)^2} \\
      &= \frac{2i}{(2i)^2} \\
      &= \frac{2i}{4i^2} \\
      &= \frac{2i}{-4} \\
      &= -\frac{i}{2}
\end{align*}

Por lo tanto, $f'(i) = -\frac{i}{2}$.

\subsection*{b) $f(z) = (z - 4i)^8$, encontrar $f'(3 + 4i)$}

Aplicando la regla de la cadena:
\begin{align*}
f'(z) &= 8(z - 4i)^7 \cdot 1 \\
      &= 8(z - 4i)^7
\end{align*}

Evaluando en $z = 3 + 4i$:
\begin{align*}
f'(3 + 4i) &= 8(3 + 4i - 4i)^7 \\
           &= 8(3)^7 \\
           &= 8 \cdot 2187 \\
           &= 17496
\end{align*}

Por lo tanto, $f'(3 + 4i) = 17496$.

\subsection*{c) $f(z) = \frac{1.5z + 2i}{3iz - 4}$, encontrar $f'(z)$ para todo $z$}

Aplicando la regla del cociente:
\begin{align*}
f'(z) &= \frac{(3iz - 4)(1.5) - (1.5z + 2i)(3i)}{(3iz - 4)^2} \\
      &= \frac{1.5(3iz - 4) - 3i(1.5z + 2i)}{(3iz - 4)^2} \\
      &= \frac{4.5iz - 6 - 4.5iz - 6i^2}{(3iz - 4)^2} \\
      &= \frac{4.5iz - 6 - 4.5iz + 6}{(3iz - 4)^2} \\
      &= \frac{0}{(3iz - 4)^2} \\
      &= 0
\end{align*}

Por lo tanto, $f'(z) = 0$ para todo $z$ tal que $3iz - 4 \neq 0$ (es decir, para $z \neq -\frac{4}{3i} = \frac{4i}{3}$).

\subsection*{d) $f(z) = i(1 - z)^n$, encontrar $f'(0)$}

Aplicando la regla de la cadena:
\begin{align*}
f'(z) &= i \cdot n(1 - z)^{n-1} \cdot (-1) \\
      &= -in(1 - z)^{n-1}
\end{align*}

Evaluando en $z = 0$:
\begin{align*}
f'(0) &= -in(1 - 0)^{n-1} \\
      &= -in(1)^{n-1} \\
      &= -in
\end{align*}

Por lo tanto, $f'(0) = -in$.

\subsection*{e) $f(z) = (iz^3 + 3z^2)^3$, encontrar $f'(2i)$}

Aplicando la regla de la cadena:
\begin{align*}
f'(z) &= 3(iz^3 + 3z^2)^2 \cdot \frac{d}{dz}(iz^3 + 3z^2) \\
      &= 3(iz^3 + 3z^2)^2 \cdot (3iz^2 + 6z) \\
      &= 3(iz^3 + 3z^2)^2(3iz^2 + 6z)
\end{align*}

Evaluando en $z = 2i$:
\begin{align*}
(2i)^2 &= 4i^2 = -4 \\
(2i)^3 &= 8i^3 = -8i \\
iz^3 + 3z^2 &= i(-8i) + 3(-4) = -8i^2 - 12 = 8 - 12 = -4 \\
3iz^2 + 6z &= 3i(-4) + 6(2i) = -12i + 12i = 0
\end{align*}

Por lo tanto:
\begin{align*}
f'(2i) &= 3(-4)^2 \cdot 0 \\
       &= 3 \cdot 16 \cdot 0 \\
       &= 0
\end{align*}

Por lo tanto, $f'(2i) = 0$.

\subsection*{f) $f(z) = \frac{z^3}{(z + i)^3}$, encontrar $f'(i)$}

Aplicando la regla del cociente:
\begin{align*}
f'(z) &= \frac{(z + i)^3 \cdot 3z^2 - z^3 \cdot 3(z + i)^2}{(z + i)^6} \\
      &= \frac{3z^2(z + i)^3 - 3z^3(z + i)^2}{(z + i)^6} \\
      &= \frac{3z^2(z + i)^2[(z + i) - z]}{(z + i)^6} \\
      &= \frac{3z^2(z + i)^2 \cdot i}{(z + i)^6} \\
      &= \frac{3iz^2}{(z + i)^4}
\end{align*}

Evaluando en $z = i$:
\begin{align*}
f'(i) &= \frac{3i(i)^2}{(i + i)^4} \\
      &= \frac{3i(-1)}{(2i)^4} \\
      &= \frac{-3i}{16i^4} \\
      &= \frac{-3i}{16} \\
      &= -\frac{3i}{16}
\end{align*}

Por lo tanto, $f'(i) = -\frac{3i}{16}$.

\newpage
\section*{Ejercicio 6 (1.6.7)}

\textbf{Demuestre que:}

\subsection*{a) $\text{Log}(-ie) = 1 - \frac{\pi}{2}i$}

Recordemos que el logaritmo complejo se define como:
$$\text{Log}(z) = \ln|z| + i\text{Arg}(z)$$

donde $\text{Arg}(z)$ es el argumento principal de $z$ con $-\pi < \text{Arg}(z) \leq \pi$.

Para $z = -ie$, tenemos:
\begin{align*}
-ie &= -i \cdot e = e(-i) = e(0 - i) = e \cdot e^{-i\pi/2}
\end{align*}

En forma cartesiana: $-ie = 0 - ie$.

Calculando el módulo:
$$|-ie| = |e(-i)| = e|i| = e \cdot 1 = e$$

Calculando el argumento:
$$\text{Arg}(-ie) = \text{Arg}(-i) = -\frac{\pi}{2}$$

Por lo tanto:
\begin{align*}
\text{Log}(-ie) &= \ln|e| + i\left(-\frac{\pi}{2}\right) \\
               &= \ln(e) - \frac{\pi}{2}i \\
               &= 1 - \frac{\pi}{2}i \quad \checkmark
\end{align*}

\subsection*{b) $\text{Log}(1 - i) = \frac{1}{2}\ln(2) - \frac{\pi}{4}i$}

Para $z = 1 - i$, tenemos:
$$|1 - i| = \sqrt{1^2 + (-1)^2} = \sqrt{2}$$

El argumento:
$$\text{Arg}(1 - i) = \arctan\left(\frac{-1}{1}\right) = -\frac{\pi}{4}$$

Por lo tanto:
\begin{align*}
\text{Log}(1 - i) &= \ln|1 - i| + i\text{Arg}(1 - i) \\
                 &= \ln(\sqrt{2}) + i\left(-\frac{\pi}{4}\right) \\
                 &= \frac{1}{2}\ln(2) - \frac{\pi}{4}i \quad \checkmark
\end{align*}

\subsection*{c) $\text{Log}(e) = 1 + 2n\pi i$}

Para $z = e$ (número real positivo), tenemos:
$$|e| = e$$

El argumento principal:
$$\text{Arg}(e) = 0$$

Sin embargo, el logaritmo complejo general (con todas las ramas) es:
$$\text{Log}(e) = \ln(e) + i(\text{Arg}(e) + 2n\pi) = 1 + i(0 + 2n\pi) = 1 + 2n\pi i$$

donde $n \in \mathbb{Z}$.

Por lo tanto:
$$\text{Log}(e) = 1 + 2n\pi i \quad \checkmark$$

\subsection*{d) $\text{Log}(i) = \left(2n + \frac{1}{2}\right)\pi i$}

Para $z = i$, tenemos:
$$|i| = 1$$

El argumento principal:
$$\text{Arg}(i) = \frac{\pi}{2}$$

El logaritmo complejo general es:
\begin{align*}
\text{Log}(i) &= \ln|i| + i(\text{Arg}(i) + 2n\pi) \\
             &= \ln(1) + i\left(\frac{\pi}{2} + 2n\pi\right) \\
             &= 0 + i\left(\frac{\pi}{2} + 2n\pi\right) \\
             &= i\left(\frac{1 + 4n}{2}\pi\right) \\
             &= \left(2n + \frac{1}{2}\right)\pi i \quad \checkmark
\end{align*}

donde $n \in \mathbb{Z}$.

\newpage
\section*{Ejercicio 10 (1.6.7)}

Las funciones hiperbólicas se definen como:
\begin{align*}
\cosh(x) &= \frac{e^x + e^{-x}}{2} \\
\sinh(x) &= \frac{e^x - e^{-x}}{2}
\end{align*}

Y las demás funciones hiperbólicas:
\begin{align*}
\tanh(x) &= \frac{\sinh(x)}{\cosh(x)} \\
\text{sech}(x) &= \frac{1}{\cosh(x)} \\
\text{csch}(x) &= \frac{1}{\sinh(x)} \\
\coth(x) &= \frac{1}{\tanh(x)}
\end{align*}

\subsection*{a) Muestre las siguientes equivalencias:}

\subsubsection*{$\cosh(x) = \cos(ix)$}

Usando la definición de $\cosh$ y la fórmula de Euler:
\begin{align*}
\cosh(x) &= \frac{e^x + e^{-x}}{2} \\
         &= \frac{e^{ix \cdot (-i)} + e^{-ix \cdot (-i)}}{2} \\
         &= \frac{e^{-ix^2} + e^{ix^2}}{2}
\end{align*}

Mejor, usando $e^{ix} = \cos(x) + i\sin(x)$ y $e^{-ix} = \cos(x) - i\sin(x)$:
\begin{align*}
\cos(ix) &= \frac{e^{i(ix)} + e^{-i(ix)}}{2} \\
        &= \frac{e^{-x} + e^{x}}{2} \\
        &= \frac{e^x + e^{-x}}{2} \\
        &= \cosh(x) \quad \checkmark
\end{align*}

\subsubsection*{$i\sinh(x) = \sin(ix)$}

\begin{align*}
\sin(ix) &= \frac{e^{i(ix)} - e^{-i(ix)}}{2i} \\
         &= \frac{e^{-x} - e^{x}}{2i} \\
         &= \frac{-(e^{x} - e^{-x})}{2i} \\
         &= \frac{e^{x} - e^{-x}}{2} \cdot \frac{-1}{i} \\
         &= \sinh(x) \cdot \frac{-i}{i^2} \\
         &= \sinh(x) \cdot \frac{-i}{-1} \\
         &= i\sinh(x) \quad \checkmark
\end{align*}

\subsubsection*{$\cos(x) = \cosh(ix)$}

\begin{align*}
\cosh(ix) &= \frac{e^{ix} + e^{-ix}}{2} \\
         &= \frac{(\cos(x) + i\sin(x)) + (\cos(x) - i\sin(x))}{2} \\
         &= \frac{2\cos(x)}{2} \\
         &= \cos(x) \quad \checkmark
\end{align*}

\subsubsection*{$i\sin(x) = \sinh(ix)$}

\begin{align*}
\sinh(ix) &= \frac{e^{ix} - e^{-ix}}{2} \\
          &= \frac{(\cos(x) + i\sin(x)) - (\cos(x) - i\sin(x))}{2} \\
          &= \frac{2i\sin(x)}{2} \\
          &= i\sin(x) \quad \checkmark
\end{align*}

\subsection*{b) Muestre las siguientes identidades:}

\subsubsection*{$\cosh^2(x) - \sinh^2(x) = 1$}

Usando las definiciones:
\begin{align*}
\cosh^2(x) - \sinh^2(x) &= \left(\frac{e^x + e^{-x}}{2}\right)^2 - \left(\frac{e^x - e^{-x}}{2}\right)^2 \\
                        &= \frac{(e^x + e^{-x})^2 - (e^x - e^{-x})^2}{4} \\
                        &= \frac{e^{2x} + 2 + e^{-2x} - (e^{2x} - 2 + e^{-2x})}{4} \\
                        &= \frac{e^{2x} + 2 + e^{-2x} - e^{2x} + 2 - e^{-2x}}{4} \\
                        &= \frac{4}{4} \\
                        &= 1 \quad \checkmark
\end{align*}

\subsubsection*{$\text{sech}^2(x) = 1 - \tanh^2(x)$}

Primero, demostremos que $\tanh^2(x) = 1 - \text{sech}^2(x)$:
\begin{align*}
\tanh^2(x) &= \frac{\sinh^2(x)}{\cosh^2(x)} \\
           &= \frac{\cosh^2(x) - 1}{\cosh^2(x)} \quad \text{(usando } \cosh^2 - \sinh^2 = 1\text{)} \\
           &= 1 - \frac{1}{\cosh^2(x)} \\
           &= 1 - \text{sech}^2(x)
\end{align*}

Por lo tanto:
$$\text{sech}^2(x) = 1 - \tanh^2(x) \quad \checkmark$$

\subsubsection*{$\cosh(2x) = \cosh^2(x) + \sinh^2(x)$}

Usando las definiciones:
\begin{align*}
\cosh(2x) &= \frac{e^{2x} + e^{-2x}}{2}
\end{align*}

Por otro lado:
\begin{align*}
\cosh^2(x) + \sinh^2(x) &= \left(\frac{e^x + e^{-x}}{2}\right)^2 + \left(\frac{e^x - e^{-x}}{2}\right)^2 \\
                        &= \frac{(e^x + e^{-x})^2 + (e^x - e^{-x})^2}{4} \\
                        &= \frac{e^{2x} + 2 + e^{-2x} + e^{2x} - 2 + e^{-2x}}{4} \\
                        &= \frac{2e^{2x} + 2e^{-2x}}{4} \\
                        &= \frac{e^{2x} + e^{-2x}}{2} \\
                        &= \cosh(2x) \quad \checkmark
\end{align*}

\subsection*{c) Resuelva las siguientes ecuaciones hiperbólicas:}

\subsubsection*{$\cosh(x) - 5\sinh(x) - 5 = 0$}

Usando las definiciones:
\begin{align*}
\cosh(x) - 5\sinh(x) - 5 &= 0 \\
\frac{e^x + e^{-x}}{2} - 5 \cdot \frac{e^x - e^{-x}}{2} - 5 &= 0 \\
\frac{e^x + e^{-x} - 5e^x + 5e^{-x}}{2} - 5 &= 0 \\
\frac{-4e^x + 6e^{-x}}{2} &= 5 \\
-4e^x + 6e^{-x} &= 10 \\
-4e^x + \frac{6}{e^x} &= 10
\end{align*}

Multiplicando por $e^x$:
\begin{align*}
-4e^{2x} + 6 &= 10e^x \\
-4e^{2x} - 10e^x + 6 &= 0 \\
4e^{2x} + 10e^x - 6 &= 0
\end{align*}

Haciendo $u = e^x$:
\begin{align*}
4u^2 + 10u - 6 &= 0 \\
2u^2 + 5u - 3 &= 0
\end{align*}

Resolviendo la ecuación cuadrática:
$$u = \frac{-5 \pm \sqrt{25 + 24}}{4} = \frac{-5 \pm 7}{4}$$

Por lo tanto:
\begin{align*}
u_1 &= \frac{2}{4} = \frac{1}{2} \quad \Rightarrow \quad e^x = \frac{1}{2} \quad \Rightarrow \quad x = \ln\left(\frac{1}{2}\right) = -\ln(2) \\
u_2 &= \frac{-12}{4} = -3 \quad \Rightarrow \quad e^x = -3 \quad \text{(no tiene solución real)}
\end{align*}

\textbf{Solución:} $x = -\ln(2)$

\subsubsection*{$2\cosh(4x) - 8\cosh(2x) + 5 = 0$}

Usando la identidad $\cosh(2x) = 2\cosh^2(x) - 1$, tenemos:
$$\cosh(4x) = 2\cosh^2(2x) - 1 = 2(2\cosh^2(x) - 1)^2 - 1$$

Pero es más simple usar directamente:
\begin{align*}
2\cosh(4x) - 8\cosh(2x) + 5 &= 0 \\
2 \cdot \frac{e^{4x} + e^{-4x}}{2} - 8 \cdot \frac{e^{2x} + e^{-2x}}{2} + 5 &= 0 \\
e^{4x} + e^{-4x} - 4(e^{2x} + e^{-2x}) + 5 &= 0
\end{align*}

Haciendo $u = e^{2x}$, entonces $e^{4x} = u^2$ y $e^{-4x} = \frac{1}{u^2}$, $e^{-2x} = \frac{1}{u}$:
\begin{align*}
u^2 + \frac{1}{u^2} - 4\left(u + \frac{1}{u}\right) + 5 &= 0 \\
u^2 + \frac{1}{u^2} - 4u - \frac{4}{u} + 5 &= 0
\end{align*}

Multiplicando por $u^2$:
\begin{align*}
u^4 + 1 - 4u^3 - 4u + 5u^2 &= 0 \\
u^4 - 4u^3 + 5u^2 - 4u + 1 &= 0
\end{align*}

Factorizando (notando que es $(u^2 - 2u + 1)^2 = (u - 1)^4$):
$$(u - 1)^4 = 0$$

Por lo tanto: $u = 1$, es decir, $e^{2x} = 1$, entonces $2x = 0$, así que $x = 0$.

\textbf{Solución:} $x = 0$

\subsubsection*{$\cosh(x) = \sinh(x) + 2\text{sech}(x)$}

\begin{align*}
\cosh(x) &= \sinh(x) + 2\text{sech}(x) \\
\cosh(x) - \sinh(x) &= \frac{2}{\cosh(x)} \\
(\cosh(x) - \sinh(x))\cosh(x) &= 2
\end{align*}

Notando que $\cosh(x) - \sinh(x) = e^{-x}$:
\begin{align*}
e^{-x}\cosh(x) &= 2 \\
e^{-x} \cdot \frac{e^x + e^{-x}}{2} &= 2 \\
\frac{1 + e^{-2x}}{2} &= 2 \\
1 + e^{-2x} &= 4 \\
e^{-2x} &= 3 \\
-2x &= \ln(3) \\
x &= -\frac{1}{2}\ln(3)
\end{align*}

\textbf{Solución:} $x = -\frac{1}{2}\ln(3)$

\end{document}

